\section{Empirical Design}
\label{sec:empirics}

\subsection{Exposure and estimating design}
\label{sec:empirical-design}

The reform has a single adoption date and no untreated group: every initial
Schedule~13D filer became subject to the five-business-day deadline on
5 February 2024 \parencite{sec2023modernization}. Identification therefore uses
cross-sectional variation in how tightly the common deadline constrains each
target, rather than a treated-versus-control comparison.

The required-trading-days measure in equation~\eqref{eq:rtd-intro} is constructed
from shares outstanding and average daily volume over the six months ending
60 trading days before the trigger. It is defined for every filing, including
one-time filers, and does not use the campaign's realised stake or outcome. The
lag reduces contamination from the activist's own trading. I will diagnose any
remaining contamination using the Item~5(c) transaction tables and repeat the
analysis with a 120-trading-day lag, following the concern in
\textcite{collindufresnefos2015} that a filer's own trading can contaminate
liquidity measures.

Under the model, the post-threshold increment is capped at
$0.05d/\mathrm{RTD}_i$. At the doubling calibration, in which the desired
increment is five percentage points, the cap binds when
$\mathrm{RTD}_i>d$. The continuous measure is the primary treatment intensity;
the binary indicator is used only to interpret that calibration.

The identifying assumption is that, absent the reform, outcomes at different
levels of required trading days would have followed parallel paths and that the
dose-response relationship would have remained stable over the estimation
window. A linear specification provides a summary:
\begin{equation}
Y_{i,q}
=\alpha_{a(i),q}+\gamma_{\ell(i),q}
+\beta\,\text{Post}_q\times\mathrm{RTD}_i+\varepsilon_{i,q},
\label{eq:dose-linear}
\end{equation}
where $\alpha_{a(i),q}$ denotes activist-by-quarter fixed effects and
$\gamma_{\ell(i),q}$ denotes liquidity-decile-by-quarter fixed effects. The
full specification also includes size-decile-by-quarter fixed effects. Since
$\beta$ is a weighted average when responses are heterogeneous, the primary
estimates are non-parametric across exposure bins:
\begin{equation}
Y_{i,q}
=\alpha_{a(i),q}+\gamma_{\ell(i),q}
+\sum_k\beta_k\mathbf{1}[\text{RTD-bin}_i=k]\times\text{Post}_q
+\varepsilon_{i,q}.
\label{eq:dose-bins}
\end{equation}
Activist-by-quarter effects absorb fund flows and changes in an activist's
strategy. Liquidity- and size-decile-by-quarter effects absorb time-varying
differences across the target distribution.

\subsection{Identification threats}
\label{sec:identification-threats}

\paragraph{Anticipation.}
The amendments were proposed on 10 March 2022, the comment period reopened on
28 April 2023, and the rule was adopted on 10 October 2023, four months before
its effective date \parencite{sec2023modernization}. The main analysis excludes
the period from October 2023 through February 2024 from the pre-reform sample.
Robustness tests also exclude windows around the March 2022 proposal and the
April 2023 reopening.

\paragraph{Bundled rule changes.}
Release 33-11253 shortened the initial Schedule~13D deadline, shortened the
deadline for Schedule~13D amendments to two business days, restructured
Schedule~13G deadlines, and required structured XML filing from 18 December
2024 \parencite{sec2023modernization}. It also declined to adopt proposed
Rule~13d-3(e) and clarified Item~6 derivative disclosure without expanding its
scope. The amendment rule
mechanically changes post-filing accumulation, so that outcome is excluded from
the behavioural tests. The primary outcomes are trigger-to-filing accumulation,
stake at filing, campaign entry and composition, and Item~6 derivative language.
I will also report switching between Schedules~13G and 13D among holders near
five per cent. A contemporaneous shift in that rate would indicate that the
Schedule~13G changes, rather than the initial Schedule~13D deadline, affect
outcomes defined over the pooled population.

\paragraph{Exposure and target characteristics.}
Required trading days depend on target size and liquidity by construction. The
SEC's Table~6 shows why realised accumulation cannot substitute for predetermined
exposure. Campaigns that had accumulated less than 90 per cent of their reported
stake by the amended deadline involved larger and more prominent activists and
larger, more liquid targets. Their average target market capitalisation was
\$1.8 billion rather than \$916 million; their Amihud ratio was 0.09 rather than
0.13; the prominent-activist share was 43.6 rather than 29.8 per cent; their
filing-window cumulative abnormal return was 17.2 rather than 5.7 per cent; and
their average increase in shareholder value was \$222 million rather than
\$36 million \parencite{sec2023modernization}. I will report covariate balance
across exposure bins before presenting estimates and include the fixed effects
described above \parencite{collindufresnefos2015,collindufresnefos2016,
edmansfangzur2013}.

An activist-level exposure measure based on past campaigns is retained only as
a heterogeneity split. It cannot identify the main effect because most filers
appear too rarely. \textcite{polk2024} report 9{,}685 initial filings by 5{,}863
entities from 2001--2022, or 1.65 filings per entity. In the project's initial
300-filing sample, 259 filer CIKs are unique, only 29 filers appear at least
twice, and only three appear at least three times. Requiring pre- and post-reform
campaigns would therefore restrict the analysis to a small group of repeat
activists that are unusually able to pre-position or use derivatives.

\subsection{Outcomes and statistical power}
\label{sec:outcomes}

The SEC's own tables imply a small population-average treatment. Only 3 per cent
of non-corporate-action campaigns, about seven per year, had accumulated less
than 90 per cent of their reported stake by the amended deadline; 1 per cent,
about one per year, had accumulated less than 75 per cent
\parencite{sec2023modernization}. Among the 20 per cent that continued to
accumulate after the deadline, only 5.9 per cent of the reported stake was
acquired within the filing window on average. Together, these figures imply an
average accumulation effect on the order of one tenth of a percentage point of
shares outstanding. The interpretation of any estimate will remain bounded by
that scale.

The expected post-reform sample contains roughly 540--700 non-corporate-action
campaigns from 2024--2026, before any repeat-activist restriction.
Table~\ref{tab:outcome-hierarchy} separates outcomes that this sample can
reasonably inform from outcomes that are descriptive.

\begin{table}[htbp]
\centering
\small
\caption{Outcome hierarchy and statistical power}
\label{tab:outcome-hierarchy}
\begin{tabular}{@{}p{0.26\linewidth}p{0.16\linewidth}p{0.48\linewidth}@{}}
\toprule
Outcome & Role & Interpretation and power \\
\midrule
Entry and target composition & primary & Best-powered extensive margin; not
selected by conditioning on a filing outcome. \\
Stake at filing and trigger-to-filing accumulation & primary & Campaign-level
tests of the deadline's accumulation effect. \\
Item~6 derivative language & primary & Campaign-level test of substitution
through the rule's treatment of cash-settled derivatives. \\
Filing delay & descriptive & Mechanical compliance outcome; not evidence of a
behavioural response. \\
Filing return and pre-filing drift & secondary & Moderately powered and dependent
on CRSP access. \\
Six- and twelve-month bid hazard & descriptive & Confidence intervals only. The
sale-intent triple interaction has a minimum detectable effect of 17--31
percentage points on a base rate of about 17 percentage points, with roughly
three bids in the materially exposed post-reform cell. \\
\bottomrule
\end{tabular}
\end{table}

\subsection{EDGAR first-stage analysis}
\label{sec:first-stage}

The first-stage analysis uses only EDGAR data and will be pre-registered against
minimum detectable effects calculated from the SEC's Tables~3, 5, and 6. It
contains five components:
\begin{enumerate}
\item untrimmed filing-delay distributions, reported separately for
corporate-action and non-corporate-action filings;
\item pre- and post-reform distributions of the stake reported at filing;
\item entry and target composition by a coarse, price-free proxy for required
trading days until CRSP-based average daily volume is available;
\item the prevalence of Item~6 cash-settled-derivative language; and
\item switching between Schedules~13G and 13D near the five-per-cent threshold.
\end{enumerate}
Filing-delay compression and accumulation after the legal deadline are
mechanical consequences of compliance and cannot establish a behavioural
response.

\subsection{Inference and falsification}
\label{sec:inference}

The principal threat is a continuous-dose estimate driven by a particular
placebo date, exposure assignment, or macroeconomic cycle. I will compare the
estimated reform effect with distributions generated from placebo reform dates
between 2015 and 2023 and from permutations of required trading days across
activists. Quarterly event studies by exposure bin begin in 2019, and a joint
pre-trend test is required before a causal interpretation. A non-US activist
block sample based on UK DTR5 and E.U.\ Transparency Directive notifications
provides a control for the global merger and interest-rate cycle. It will not be
used as a legal-threshold comparison because the disclosure regimes are not
otherwise comparable.

\subsection{The timing of target gains}
\label{sec:gain-anatomy}

Total target gain is divided into the filing return
$G_{\mathrm{fil}}=\log(P^F/P^0)$, the post-filing run-up
$G_{\mathrm{run}}=\log(P^B/P^F)$, and the bid markup
$G_{\mathrm{mkp}}=\log(B/P^B)$. The unaffected price $P^0$ is measured
60 trading days before the trigger in both reform periods, with robustness at
30 and 120 days. Fixing this anchor prevents the treatment from moving the
starting point of the decomposition.

The reform shortened the median trigger-to-filing interval from seven to five
business days in the project's initial data. This mechanically shortens the
calendar span in the filing return and lengthens the span in the run-up, even if
behaviour does not change. The primary filing measure is therefore the
filing-date $[-1,+1]$ cumulative abnormal return as a share of total gain, with
the trigger-to-bid horizon held fixed. A mechanical placebo recuts pre-reform
campaigns as though they had filed under a five-business-day rule. A behavioural
placebo restricts the sample to campaigns that already filed within five
business days under the old rule, for which the reform did not change the legal
window in practice.

If an M\&A outcomes database is unavailable, I will hand-collect the
acquired-target subsample from Forms~8-K and DEFM14A using standard conventions
for toeholds, premia, and deal outcomes
\parencite{bett2014,eckbomalenkothorburn2025}. Cross-sectional tests use industry
M\&A intensity as a proxy for fringe-bidder arrival, the strategic-buyer share
as a proxy for the portability of activist improvements, and a discontinuity
check around the five-per-cent ownership threshold
\parencite{corumlevit2019,burkartlee2022,brav2008}. The capitalisation
interpretation remains a measurement restriction, not a separate contribution
\parencite{celentanolevine2025}.

\subsection{Data feasibility}
\label{sec:data-feasibility}

The core sample spans 2022--2026. This window brackets the reform while limiting
the Item~5(c) extraction burden; the SEC's published 2011--2021 tables provide a
pre-period benchmark. EDGAR supports the first-stage analysis. CRSP or WRDS
access is required for average daily volume and return-based outcomes, and an
M\&A source is required for the full gain decomposition. These access decisions
govern the price-based extensions, not the EDGAR analysis. The companion
execution plan records the acquisition schedule, validation thresholds, and
fallbacks.
